\documentclass[10pt,a4paper]{amsart}
\usepackage[margin=19mm]{geometry}
\usepackage{amsmath,amssymb,mathtools}
\usepackage[scr=rsfs,cal=euler]{mathalfa}
\usepackage{xfrac}
\usepackage[unicode,hidelinks,bookmarks=false]{hyperref}
\newcommand{\ie}{i.e.\ }
\newcommand{\cf}{cf.\ }
\newcommand{\resp}{respectively\ }
\newcommand{\Subsection}{\subsection}
\providecommand{\nobreakdash}{\nobreak\hbox{-}}
\setlength{\emergencystretch}{2em}
\multlinegap0pt
\usepackage{bm}
\usepackage{bbm}
\usepackage{dsfont}
\usepackage{xspace}
\usepackage{cancel}
\usepackage{mathtools}
\usepackage{amsthm}
\makeatletter
\@ifundefined{lemma}{\newtheorem{lemma}{Lemma}}{}
\makeatother
\title[Sharp power concavity of two relevant free boundary problems]{Sharp power concavity of two relevant free boundary problems of reaction-diffusion type}
\author{Qingyou He}
\date{Source: arXiv:2509.08768v1 (CC BY 4.0)}
\begin{document}
\maketitle

\noindent\emph{Source and licence.} Sharp power concavity of two relevant free boundary problems of reaction-diffusion type. Version arXiv:2509.08768v1 (CC BY 4.0), distributed under the Creative Commons Attribution 4.0 International licence, as recorded in the arXiv licence metadata for this exact version and shown on the abstract page https://arxiv.org/abs/2509.08768v1. Full rights notice: licenses/arxiv-2509.08768-CC-BY-4.0.txt.\par\smallskip

\noindent\textit{Editorial scope.} Counted unit: the Lemma whose source label is \texttt{npcl} in the article (source0.tex lines 490--497, source label \texttt{npcl}), delivered together with the article's own complete original proof of it (source0.tex lines 498--613, one \texttt{proof} environment of 116 lines). No mathematics is written, repaired or re-derived: statement and proof are the article's own text line for line. Required context retained uncounted: g0b (lines 187--189); partialt11 (lines 466--474); partialt110 (lines 476--482). Every definition, display and label of those lines is retained unchanged. Omitted from this package and not redistributed: 1--186, 190--465, 475--475, 483--489, 614--1321. Nothing inside a retained excerpt is omitted or altered: every retained body is delivered exactly as the article's lines are written, so no omission receipt of any kind accompanies this package. Citations: the retained excerpts carry no citation command at all, so no bibliography entry of the article is transcribed and no reference is redistributed. Reference adaptation: none was needed and none was made; every reference inside the delivered unit resolves inside it, so no unresolved reference or citation is produced. Printed slips kept literal: no printed word, symbol, hypothesis or number of the counted statement or of its proof was altered. Layout and class: the article's own class is \texttt{article} with packages including mathrsfs, amsmath, amsfonts, amssymb, amsthm, caption, subfigure, cite, color, graphicx, float, paralist, hyperref, indentfirst; the wrapper is the lane's pinned \texttt{amsart} editorial wrapper at 10pt on A4, which restates the theorem-like environments the retained text needs and adds the article's own macro definitions that the retained text needs (none), with extra wrapper packages (bm, bbm, dsfont, xspace, cancel, mathtools). The delivered numbering is the unit's own. Assets: no retained span contains a figure environment, a graphics-inclusion command, a PDF-page inclusion, a table or a TikZ picture; no image file is copied into this package, the package declares no asset dependency and no image file is redistributed with it. Licence: version arXiv:2509.08768v1 of this article is distributed under the Creative Commons Attribution 4.0 International licence, as recorded in the arXiv licence metadata for this exact version and shown on the abstract page https://arxiv.org/abs/2509.08768v1. A full rights notice is delivered at licenses/arxiv-2509.08768-CC-BY-4.0.txt. Characters: every retained excerpt is pure ASCII, so the delivered engine, XeLaTeX with its default Latin Modern text font, reports no missing character.
\begin{equation}\label{g0b}
|G'(P)|+|PG''(P)|\leq AP^{\gamma-1}\quad \text{for }P\in(0,P_H] \text{ and some }A, \gamma>0,
\end{equation}

\begin{equation}\label{partialt11}
\begin{aligned}
    \partial_t w_{11}=&rw^{\frac{1}{\alpha}}w_{kk11}+\frac{2r}{\alpha}w^{\frac{1}{\alpha}-1}w_1w_{kk1}+\frac{r}{\alpha}(\frac{1}{\alpha}-1)w^{\frac{1}{\alpha}-2}w_1^2w_{kk}+\frac{r}{\alpha}w^{\frac{1}{\alpha}-1}w_{11}w_{kk}\\
    &+\frac{1}{\alpha}\big(1+r(1-\alpha)\big)\{(\frac{1}{\alpha}-2)(\frac{1}{\alpha}-1)w^{\frac{1}{\alpha}-3}w_1^2w_k^2\\
    &+(\frac{1}{\alpha}-1)w^{\frac{1}{\alpha}-2}w_{11}w_k^2+4(\frac{1}{\alpha}-1)w^{\frac{1}{\alpha}-3}w_1w_kw_{k1}\\
    &+2w^{\frac{1}{\alpha}-1}w_{1k}^2+2w^{\frac{1}{\alpha}-1}w_kw_{k11}\}+rw_{11}\big(\alpha G(w^{\frac{1}{\alpha}})+w^{\frac{1}{\alpha}}G'(w^{\frac{1}{\alpha}})\big)\\
    &+r|w_1|^2\big[(\frac{1}{\alpha}+1)w^{\frac{1}{\alpha}-1}G'(w^{\frac{1}{\alpha}})+\frac{1}{\alpha}w^{\frac{2}{\alpha}-1}G''(w^{\frac{1}{\alpha}})\big],
\end{aligned}
\end{equation}

\begin{equation}\label{partialt110}
\begin{aligned}
\partial_t w_{11}=&re^ww_{kk11}+2re^{w}w_1w_{kk1}+re^ww_1^2w_{kk}+re^ww_{11}w_{kk}+me^ww_1^2w_k^2\\
&+me^ww_1w_kw_{k1}+2me^ww_{1k}^2+2me^ww_kw_{k11}\\
&+rw_{1}^2\big(G''(e^w)e^{2w}+G'(e^w)e^w\big)+rw_{11} G'(e^w)e^w.
\end{aligned}    
\end{equation}

\begin{lemma}\label{npcl}
 For each $\alpha\in[0,1]\backslash\{\frac{1}{2}\}$, there exists a $r>0$  and a smooth  positive function $w$ on $\overline{B_r(0)}$ such that 
 \begin{itemize}
    \item[(\expandafter{\romannumeral1})] $w_{11}(0)=0$, $w_{ii}(0)<0$ for $2\leq i\leq N$ and $w_{ij}(0)=0$ otherwise.
    \item[(\expandafter{\romannumeral2})] $\big(D^{2}w\big)<0$ on $\overline{B_{\rho}(0)}\backslash \{0\}$.
    \item[(\expandafter{\romannumeral3})] $\partial_t w_{11}>0$ from the right hand side \eqref{partialt11}-\eqref{partialt110} at the origin $(0,0)$.
\end{itemize}
\end{lemma}

\begin{proof}
The proof is divided into two cases. 

\noindent\textbf{\emph{Case 1, $\alpha\in[0,\frac{1}{2})$.}} Define 
\begin{equation}\label{w}w=c+x_1-x_1^4+\sum\limits_{i=2}^{N}(ax_1x_i^2-x_i^2-2a^2x_1^2x_i^2\big),\end{equation}
where the constants $a,c>0$ are determined later. By direct computations, we get 
\begin{align}
 w_{11}=&-12x_1^2-\sum\limits_{i=2}^{N}4a^2x_i^2,\notag\\
 w_{1i}=&w_{i1}=2ax_i-8a^2x_1x_i\quad \text{ for }i=2,...,N,\notag\\
 w_{ii}=&2a x_1-2-4x_1^2\quad \text{ for }i=2,...,N,\notag\\
 w_{ij}=&0\quad \text{for }i,j=2,...,N,\notag
\end{align}
Then, (\expandafter{\romannumeral1}) naturally holds.

\vspace{2mm}

For (\expandafter{\romannumeral2}) on the Hessian matrix $(D^2w)$, we calculate directly determinant of the $j$-$th$ leading principle minor $(D^2w)_j$ for $j=2,...,N$, that is   
\[\text{determinant of } (D^2w)_j=(-2)^{j}(6x_1^2+2a^2\sum\limits_{i=2}^Nx_i^2-a^2\sum\limits_{i=2}^{j}x_i^2)+\text{higher order terms}.\]
Due to
\[6x_1^2+2a^2\sum\limits_{i=2}^Nx_i^2-a^2\sum\limits_{i=2}^{j}x_i^2>0 \text{ for }x\neq0, \]
it holds by taking small $\rho>0$ that 
\[(-1)^j\big(\text{determinant of } (D^2w)_j\big)>0,\quad x\in \overline{B_{\rho}}\backslash\{0\}.\]
Hence, $(D^2w)$ is a strictly negative-definite matrix in $\overline{B_{\rho}}\backslash\{0\}$, so (\expandafter{\romannumeral2}) holds.


\vspace{2mm}

We turn to (\expandafter{\romannumeral3}), at the origin $(0,0)$, it follows from direct calculations that  
\begin{align}
&w=c,         &&w_{1}=1,         &&&w_{kk}=-2,\notag\\
&w_{kk1}=2a,  &&w_{1111}=-24,      &&&w_{kk11}=-8a^2.\notag
\end{align}
Then,  the right hand side of \eqref{partialt11} at the origin $(0,0)$ is 
\begin{align}R(a,c)=&-rc\big(24+8(N-1)\big)a^2+\frac{2r}{\alpha}c^{\frac{1}{\alpha}-1}a+\frac{1}{\alpha}\big(1+r(1-\alpha)\big)(\frac{1}{\alpha}-2)(\frac{1}{\alpha}-1)a^4\notag\\
&+r\big[(\frac{1}{\alpha}+1)c^{\frac{1}{\alpha}-1}G'(c^{\frac{1}{\alpha}})+\frac{1}{\alpha}c^{\frac{2}{\alpha}-1}G''(c^{\frac{1}{\alpha}})\big].\notag\end{align}
For $\alpha\in (0,\frac{1}{2})$, it holds by $a\gg1$, and $|G'(c^{\frac{1}{\alpha}})|+|G''(c^{\frac{1}{\alpha}})|<\infty$ for some $c>0$ that 
\begin{equation}\label{012}
\partial_t w_{11}(0,0)=R(a,c)>0.
\end{equation}

 In the case of $\alpha=0$, let $|G'(e^c)|+|G''(e^c)|<\infty$ for some $c>0$,  the right hand side of \eqref{partialt110} with sufficiently large $a>0$ yields
\begin{equation}\label{0}
\begin{aligned}
\partial_t w_{11}(0,0)=&-r\big(24+8(N-1)\big)\partial_t w_{11}(0,0)+4r(N-1)a-2r(N-1)a^2+ma^4\\
&+ra^2\big(G''(e^c)e^{2c}+G'(e^c)e^c\big)>0.
\end{aligned}
 \end{equation}
Hence,  (\expandafter{\romannumeral3}) is justified based on the above computations \eqref{012}-\eqref{0} for $\alpha\in[0,\frac{1}{2})$.

\noindent\textbf{\emph{Case 2, $\alpha=1$.}} Still for the function $w$  given by \eqref{w} with $\alpha=1$, (\expandafter{\romannumeral1}) and (\expandafter{\romannumeral2}) naturally hold as \textbf{\emph{Case 1.}}
By means of he condition \eqref{g0b} on $G$, we have 
\[\begin{aligned}
\partial_t w_{11}(0,0)=&2a+[2G'(c)+cG''(c)]-ca^2\big(24+8(N-1)\big)\\
\geq &2a-ca^2\big(24+8(N-1)\big)-Ac^{\gamma-1}.
\end{aligned}\]
 For $\gamma\geq 1$, let $c:=a^{-2}$, we take large $a>0$ such that 
\begin{equation}\label{ptw111}\partial_t w_{11}(0,0)\geq 2a-24-8(N-1)-Aa^{-2(\gamma-1)}>0.\end{equation}
For $\gamma\in(0,1)$, it holds by setting $c=:(Aa)^{\frac{-1}{1-\gamma}}$ that 
\begin{equation}\label{ptw112}\partial_t w_{11}(0,0)\geq2a-a^{2-\frac{1}{1-\gamma}}(A)^{\frac{-2}{1-\gamma}}\big(24+8(N-1)\big)-a>0 \end{equation}
for large $a>0$. (\expandafter{\romannumeral3}) is valid for $\alpha=1$ on account of \eqref{ptw111}-\eqref{ptw112}.















\vspace{2mm}

\noindent\textbf{\emph{Case 3, $\alpha\in(\frac{1}{2},1)$.}}     Redefine 
\[w=c+\frac{\alpha(\frac{3}{2}-\alpha)^{\frac{1}{2}}}{b(1-\alpha)}x_1-\frac{x_1^4}{12b^2}+\sum\limits_{i=2}^{N}\big(-b^2x_i^2+b(\frac{3}{2}-\alpha)^{\frac{1}{2}}x_1x_i^2-x_1^2x_i^2\big),\]
where $b>0$ is to be determined later. The each component of the Hessian matrix $(D^2w)$ is 
\begin{align}
w_{11}&=-\frac{x_1^2}{b^2}-2\sum\limits_{i=2}^{N}x_i^2\notag\\
w_{1i}&=w_{i1}=2b(\frac{3}{2}-\alpha)^{\frac{1}{2}}x_i-4x_1x_i\quad\text{ for }i=2,...,N,\notag\\
w_{ii}&=-2b^2+2b(\frac{3}{2}-\alpha)^{\frac{1}{2}}x_1-2x_1^2\quad\text{ for }i=2,...,N,\notag\\
w_{ij}&=0\quad \text{for }i,j=2,...,N,\notag
\end{align}
which verifies (\expandafter{\romannumeral1}) by taking the value at the origin.

\vspace{2mm}

Similar to the computations of \textbf{\emph{Case 1}}, the determinant of the $j$-th leading principal minor is given by 
\[\text{determinant of } (D^2w)_j=(-1)^{j}(2b^2)^{j-1}(\frac{x_1^2}{b^2}+2\sum\limits_{i=2}^Nx_i^2-2(\frac{3}{2}-\alpha)\sum\limits_{i=2}^{j}x_i^2)+\text{higher order terms}.\]
Since 
\[\frac{x_1^2}{b^2}+2\sum\limits_{i=2}^Nx_i^2-2(\frac{3}{2}-\alpha)\sum\limits_{i=2}^{j}x_i^2>0,\quad\alpha\in (\frac{1}{2},1),\quad x\neq0,\]
it concludes by taking small $\rho>0$ that 
\[(-1)^j\big(\text{determinant of } (D^2w)_j\big)>0,\quad x\neq0,\ |x|\leq \rho,\]
which proves (\expandafter{\romannumeral2}).

\vspace{2mm}

We are going to verify (\expandafter{\romannumeral3}).  Due to the explicit form of $w$,  it is direct to obtain  
\begin{align}
&w=c,         &&w_{1}=\frac{\alpha(\frac{3}{2}-\alpha)^{\frac{1}{2}}}{b(1-\alpha)},         &&&w_{kk}=-2b^2,\notag\\
&w_{kk1}=2b(\frac{3}{2}-\alpha)^{\frac{1}{2}},  &&w_{1111}=-\frac{2}{b^2},     &&&\ \ \ w_{kk11}=-4.\notag
\end{align}
Substituting these values into \eqref{partialt11} at the origin, for $|G'(c)|+|G''(c)|<\infty$ with some $c>0$, we attain
\begin{align}
\partial_t w_{11}=&r(-\frac{2}{b^2}-4(N-1))+2r\frac{\frac{3}{2}-\alpha}{1-\alpha}(N-1)\notag+\frac{1}{\alpha}\big(1+r(1-\alpha)\big)(\frac{1}{\alpha}-2)(\frac{1}{\alpha}-1)\big(\frac{\alpha^2(\frac{3}{2}-\alpha)}{b^2(1-\alpha)^2}\big)^{2}\notag\\
&+r\frac{\alpha^2(\frac{3}{2}-\alpha)}{b^2(1-\alpha)^2}\big[(\frac{1}{\alpha}+1)G'(c)+\frac{1}{\alpha}G''(c)\big]\notag\\
=&\frac{\mathcal{A}_{\alpha,m}}{b^2}-\frac{\mathcal{B}_{\alpha,m}}{b^4}+r(N-1)\big(\frac{2\alpha-1}{1-\alpha}\big)\notag
\end{align}
with \begin{align}\mathcal{A}_{\alpha,m}&:=-2r+r\frac{\alpha^2(\frac{3}{2}-\alpha)}{(1-\alpha)^2}\big[(\frac{1}{\alpha}+1)G'(c)+\frac{1}{\alpha}G''(c)\big],\notag\\
\mathcal{B}_{\alpha,m}&:=\frac{1}{\alpha}\big(1+r(1-\alpha)\big)(2-\frac{1}{\alpha})(\frac{1}{\alpha}-1)\big(\frac{\alpha^2(\frac{3}{2}-\alpha)}{(1-\alpha)^2}\big)^{2}.\notag\end{align}
Therefore, due to $r(N-1)\big(\frac{2\alpha-1}{1-\alpha})>0$, 
choosing large $b>0$ means that  (\expandafter{\romannumeral3}) holds.
\end{proof}

\end{document}
